% 来源及取材范围见分题文件；此为整理者转录的正文，不是下载原件。
\paragraph{The group $\mu_n$.}
Let $n\in\mathbf N$ and consider the functor
\[
\mathit{Sch}^{opp}\longrightarrow\mathit{Ab},\qquad S\longmapsto\mu_n(S)=\{t\in\Gamma(S,\mathcal O_S^*)\mid t^n=1\}.
\]
By Groupoids, Example 39.5.2 this is a representable functor, and the scheme representing it is denoted $\mu_n$ also. By Lemma 59.15.8 this functor satisfies the sheaf condition for the fpqc topology (in particular, it also satisfies the sheaf condition for the étale, Zariski, etc topology).

\begin{lemma}[59.28.1]
If $n\in\mathcal O_S^*$ then
\[0\to\mu_{n,S}\to\mathbf G_{m,S}\xrightarrow{(\cdot)^n}\mathbf G_{m,S}\to0\]
is a short exact sequence of sheaves on both the small and big étale site of $S$.
\end{lemma}
\begin{proof}
By definition the sheaf $\mu_{n,S}$ is the kernel of the map $(\cdot)^n$. Hence it suffices to show that the last map is surjective. Let $U$ be a scheme over $S$. Let $f\in\mathbf G_m(U)=\Gamma(U,\mathcal O_U^*)$. We need to show that we can find an étale cover of $U$ over the members of which the restriction of $f$ is an $n$th power. Set
\[U'=\underline{\operatorname{Spec}}_U(\mathcal O_U[T]/(T^n-f))\xrightarrow{\pi}U.\]
(See Constructions, Section 27.3 or 27.4 for a discussion of the relative spectrum.) Let $\operatorname{Spec}(A)\subset U$ be an affine open, and say $f|_{\operatorname{Spec}(A)}$ corresponds to the unit $a\in A^*$. Then $\pi^{-1}(\operatorname{Spec}(A))=\operatorname{Spec}(B)$ with $B=A[T]/(T^n-a)$.

The ring map $A\to B$ is finite free of rank $n$, hence it is faithfully flat, and hence we conclude that $\operatorname{Spec}(B)\to\operatorname{Spec}(A)$ is surjective. Since this holds for every affine open in $U$ we conclude that $\pi$ is surjective. In addition, $n$ and $T^{n-1}$ are invertible in $B$, so $nT^{n-1}\in B^*$ and the ring map $A\to B$ is standard étale, in particular étale. Since this holds for every affine open of $U$ we conclude that $\pi$ is étale. Hence $\mathcal U=\{\pi:U'\to U\}$ is an étale covering. Moreover, $f|_{U'}=(f')^n$ where $f'$ is the class of $T$ in $\Gamma(U',\mathcal O_{U'}^*)$, so $\mathcal U$ has the desired property.
\end{proof}

\begin{lemma}[59.28.3]
For any $n\in\mathbf N$ the sequence
\[0\to\mu_{n,S}\to\mathbf G_{m,S}\xrightarrow{(\cdot)^n}\mathbf G_{m,S}\to0\]
is a short exact sequence of sheaves on the site $(\mathit{Sch}/S)_{fppf}$ and $(\mathit{Sch}/S)_{syntomic}$.
\end{lemma}
\begin{proof}
By definition the sheaf $\mu_{n,S}$ is the kernel of the map $(\cdot)^n$. Hence it suffices to show that the last map is surjective. Since the syntomic topology is weaker than the fppf topology, see Topologies, Lemma 34.7.2, it suffices to prove this for the syntomic topology. Let $U$ be a scheme over $S$. Let $f\in\mathbf G_m(U)=\Gamma(U,\mathcal O_U^*)$. We need to show that we can find a syntomic cover of $U$ over the members of which the restriction of $f$ is an $n$th power. Set
\[U'=\underline{\operatorname{Spec}}_U(\mathcal O_U[T]/(T^n-f))\xrightarrow{\pi}U.\]
(See Constructions, Section 27.3 or 27.4 for a discussion of the relative spectrum.) Let $\operatorname{Spec}(A)\subset U$ be an affine open, and say $f|_{\operatorname{Spec}(A)}$ corresponds to the unit $a\in A^*$. Then $\pi^{-1}(\operatorname{Spec}(A))=\operatorname{Spec}(B)$ with $B=A[T]/(T^n-a)$.

The ring map $A\to B$ is finite free of rank $n$, hence it is faithfully flat, and hence we conclude that $\operatorname{Spec}(B)\to\operatorname{Spec}(A)$ is surjective. Since this holds for every affine open in $U$ we conclude that $\pi$ is surjective. In addition, $B$ is a global relative complete intersection over $A$, so the ring map $A\to B$ is standard syntomic, in particular syntomic. Since this holds for every affine open of $U$ we conclude that $\pi$ is syntomic.

Hence $\mathcal U=\{\pi:U'\to U\}$ is a syntomic covering. Moreover, $f|_{U'}=(f')^n$ where $f'$ is the class of $T$ in $\Gamma(U',\mathcal O_{U'}^*)$, so $\mathcal U$ has the desired property.
\end{proof}

\paragraph{Pairs and their group operation (59.28.4).}
Let $n\in\mathbf N$ and let $S$ be any scheme. There is another more direct way to describe the first cohomology group with values in $\mu_n$. Consider pairs $(\mathcal L,\alpha)$ where $\mathcal L$ is an invertible sheaf on $S$ and $\alpha:\mathcal L^{\otimes n}\to\mathcal O_S$ is a trivialization of the $n$th tensor power of $\mathcal L$. Let $(\mathcal L',\alpha')$ be a second such pair. An isomorphism $\varphi:(\mathcal L,\alpha)\to(\mathcal L',\alpha')$ is an isomorphism $\varphi:\mathcal L\to\mathcal L'$ of invertible sheaves such that the diagram
\[
\xymatrix{\mathcal L^{\otimes n}\ar[d]_{\varphi^{\otimes n}}\ar[r]_\alpha&\mathcal O_S\ar[d]^1\\(\mathcal L')^{\otimes n}\ar[r]^{\alpha'}&\mathcal O_S}
\]
commutes. Thus we have
\[
\mathit{Isom}_S((\mathcal L,\alpha),(\mathcal L',\alpha'))=
\begin{cases}\emptyset&\text{if they are not isomorphic},\\
H^0(S,\mu_{n,S})\cdot\varphi&\text{if $\varphi$ is an isomorphism of pairs}.
\end{cases}\tag{59.28.4.1}
\]
Moreover, given two pairs $(\mathcal L,\alpha)$, $(\mathcal L',\alpha')$ the tensor product
\[(\mathcal L,\alpha)\otimes(\mathcal L',\alpha')=(\mathcal L\otimes\mathcal L',\alpha\otimes\alpha')\]
is another pair. The pair $(\mathcal O_S,1)$ is an identity for this tensor product operation, and an inverse is given by
\[(\mathcal L,\alpha)^{-1}=(\mathcal L^{\otimes-1},\alpha^{\otimes-1}).\]
Hence the collection of isomorphism classes of pairs forms an abelian group. Note that
\[(\mathcal L,\alpha)^{\otimes n}=(\mathcal L^{\otimes n},\alpha^{\otimes n})\xrightarrow{\alpha}(\mathcal O_S,1)\]
is an isomorphism hence every element of this group has order dividing $n$. We warn the reader that this group is in general not the $n$-torsion in $\operatorname{Pic}(S)$.

\begin{lemma}[59.28.5, Tag 040Q]
Let $S$ be a scheme. There is a canonical identification
\[H^1_{\acute etale}(S,\mu_n)=\text{group of pairs $(\mathcal L,\alpha)$ up to isomorphism as above}\]
if $n$ is invertible on $S$. In general we have
\[H^1_{fppf}(S,\mu_n)=\text{group of pairs $(\mathcal L,\alpha)$ up to isomorphism as above}.\]
The same result holds with fppf replaced by syntomic.
\end{lemma}
\begin{proof}
We first prove the second isomorphism. Let $(\mathcal L,\alpha)$ be a pair as above. Choose an affine open covering $S=\bigcup U_i$ such that $\mathcal L|_{U_i}\cong\mathcal O_{U_i}$. Say $s_i\in\mathcal L(U_i)$ is a generator. Then $\alpha(s_i^{\otimes n})=f_i\in\mathcal O_S^*(U_i)$.

Writing $U_i=\operatorname{Spec}(A_i)$ we see there exists a global relative complete intersection $A_i\to B_i=A_i[T]/(T^n-f_i)$ such that $f_i$ maps to an $n$th power in $B_i$. In other words, setting $V_i=\operatorname{Spec}(B_i)$ we obtain a syntomic covering $\mathcal V=\{V_i\to S\}_{i\in I}$ and trivializations $\varphi_i:(\mathcal L,\alpha)|_{V_i}\to(\mathcal O_{V_i},1)$.
We will use this result (the existence of the covering $\mathcal V$) to associate to this pair a cohomology class in $H^1_{syntomic}(S,\mu_{n,S})$. We give two (equivalent) constructions.

First construction: using \v Cech cohomology. Over the double overlaps $V_i\times_S V_j$ we have the isomorphism
\begin{gather*}
(\mathcal O_{V_i\times_S V_j},1)
\xrightarrow{\operatorname{pr}_0^*\varphi_i^{-1}}
(\mathcal L|_{V_i\times_S V_j},\alpha|_{V_i\times_S V_j})\\
\xrightarrow{\operatorname{pr}_1^*\varphi_j}(\mathcal O_{V_i\times_S V_j},1)
\end{gather*}
of pairs. By (59.28.4.1) this is given by an element $\zeta_{ij}\in\mu_n(V_i\times_S V_j)$. We omit the verification that these $\zeta_{ij}$'s give a $1$-cocycle, i.e., give an element $(\zeta_{i_0i_1})\in\check C(\mathcal V,\mu_n)$ with $d(\zeta_{i_0i_1})=0$. Thus its class is an element in $\check H^1(\mathcal V,\mu_n)$ and by Theorem 59.19.2 it maps to a cohomology class in $H^1_{syntomic}(S,\mu_{n,S})$.

Second construction: Using torsors. Consider the presheaf
\[\mu_n(\mathcal L,\alpha):U\longmapsto\mathit{Isom}_U((\mathcal O_U,1),(\mathcal L,\alpha)|_U)\]
on $(\mathit{Sch}/S)_{syntomic}$. We may view this as a subpresheaf of $\mathcal Hom_{\mathcal O}(\mathcal O,\mathcal L)$ (internal hom sheaf, see Modules on Sites, Section 18.27). Since the conditions defining this subpresheaf are local, we see that it is a sheaf. By (59.28.4.1) this sheaf has a free action of the sheaf $\mu_{n,S}$. Hence the only thing we have to check is that it locally has sections. This is true because of the existence of the trivializing cover $\mathcal V$. Hence $\mu_n(\mathcal L,\alpha)$ is a $\mu_{n,S}$-torsor and by Cohomology on Sites, Lemma 21.4.3 we obtain a corresponding element of $H^1_{syntomic}(S,\mu_{n,S})$.

Ok, now we have to still show the following
\begin{enumerate}
\item The two constructions give the same cohomology class.
\item Isomorphic pairs give rise to the same cohomology class.
\item The cohomology class of $(\mathcal L,\alpha)\otimes(\mathcal L',\alpha')$ is the sum of the cohomology classes of $(\mathcal L,\alpha)$ and $(\mathcal L',\alpha')$.
\item If the cohomology class is trivial, then the pair is trivial.
\item Any element of $H^1_{syntomic}(S,\mu_{n,S})$ is the cohomology class of a pair.
\end{enumerate}
We omit the proof of (1). Part (2) is clear from the second construction, since isomorphic torsors give the same cohomology classes. Part (3) is clear from the first construction, since the resulting \v Cech classes add up. Part (4) is clear from the second construction since a torsor is trivial if and only if it has a global section, see Cohomology on Sites, Lemma 21.4.2.

Part (5) can be seen as follows (although a direct proof would be preferable). Suppose $\xi\in H^1_{syntomic}(S,\mu_{n,S})$. Then $\xi$ maps to an element $\overline\xi\in H^1_{syntomic}(S,\mathbf G_{m,S})$ with $n\overline\xi=0$. By Theorem 59.24.1 we see that $\overline\xi$ corresponds to an invertible sheaf $\mathcal L$ whose $n$th tensor power is isomorphic to $\mathcal O_S$. Hence there exists a pair $(\mathcal L,\alpha')$ whose cohomology class $\xi'$ has the same image $\overline{\xi'}$ in $H^1_{syntomic}(S,\mathbf G_{m,S})$. Thus it suffices to show that $\xi-\xi'$ is the class of a pair. By construction, and the long exact cohomology sequence above, we see that $\xi-\xi'=\partial(f)$ for some $f\in H^0(S,\mathcal O_S^*)$. Consider the pair $(\mathcal O_S,f)$. We omit the verification that the cohomology class of this pair is $\partial(f)$, which finishes the proof of the first identification (with fppf replaced with syntomic).

To see the first, note that if $n$ is invertible on $S$, then the covering $\mathcal V$ constructed in the first part of the proof is actually an étale covering (compare with the proof of Lemma 59.28.1). The rest of the proof is independent of the topology, apart from the very last argument which uses that the Kummer sequence is exact, i.e., uses Lemma 59.28.1.
\end{proof}
